
\documentclass[mn,a4paper,fleqn]{w-art}
\usepackage{times,cite,w-thm}
\usepackage{amssymb,amsmath,enumerate,graphics,graphicx} 


\textwidth=16cm  %
\textheight=21cm  %
\oddsidemargin=-0.1cm  %
\evensidemargin=-0.1cm  %
\newcommand{\meas}{{\rm meas}\,}
\newcommand{\diam}{{\rm diam}\,}
\newcommand{\dist}{{\rm dist}\,}
\newcommand{\Area}{{\rm area}\,}
\newcommand{\CAP}{{\rm cap}\,}
\newcommand{\be}{\begin{equation}}
\newcommand{\ee}{\end{equation}}
\newcommand{\bl}{\begin{lemma}}
\newcommand{\el}{\end{lemma}}


%\newtheorem{theorem}{Theorem}
%\newtheorem{lemma}{Lemma}
%\newtheorem{corollary}{Corollary}
\newtheorem{theo}{Theorem}
%\newtheorem{conjecture}{Conjecture}
\newtheorem{lemmaA}{Lemma}
\newcommand{\blemmaA}{\begin{lemmaA}}
\newcommand{\elemmaA}{\end{lemmaA}}
%\newtheorem{proposition}{Proposition}
\newcommand{\bprop}{\begin{proposition}}
\newcommand{\eprop}{\end{proposition}}




\newcommand{\bt}{\begin{theorem}}
\newcommand{\et}{\end{theorem}}
\newcommand{\bc}{\begin{corollary}}
\newcommand{\ec}{\end{corollary}}
\newcommand{\bcon}{\begin{conjecture}}
\newcommand{\econ}{\end{conjecture}}
\newcommand{\ep}{\varepsilon}
\newcommand{\ph}{\varphi}
\newcommand{\va}{\varphi}
\newcommand{\cA}{{\cal{A}}}
\newcommand{\al}{\alpha}
%\newcommand{\th}{\theta}
\newcommand{\ga}{\gamma}
\newcommand{\Res}{{\rm{Res}}}
\newcommand{\CU}{{\overline{\U}}}
\newcommand{\U}{{\mathbb U}}
\newcommand{\II}{\int\!\!\int}
\newcommand{\D}{{\cal D}}
\newcommand{\DD}{{\mathbb{D}}}
\newcommand{\e}{\varepsilon}
\newcommand{\Om}{\Omega}
\newcommand{\UH}{{\mathbb {H}}}

\newcommand{\Ha}{{\cal H}}
\newcommand{\C}{{\mathbb C}}
\newcommand{\CC}{{\overline{\C}}}
\newcommand{\R}{{\mathbb R}}
\newcommand{\T}{{\mathbb T}}
\newcommand{\A}{{\cal A}}
\newcommand{\F}{{\cal F}}
\newcommand{\CR}{{\overline{\R}}}

\newtheorem{fig}{Box for Figure}
\renewcommand{\thetheo}{\Alph{theo}}
\renewcommand{\theequation}{\thesection.\arabic{equation}}
\renewcommand{\thelemmaA}{\Alph{lemmaA}}

\newcommand{\pf}{\textit{Proof. \@}}                  % Proof
\newcommand{\epf}{\ensuremath{\Box}}                  % End proof
%%%%%%%%%%%%%%%%%%%%%%%%%%%%%%%%%%%%%%%%%%%%%%%%%%%%%%%%%%%%%%%%%%%%%%%%%%%%%%%%%%%
\newcommand{\Na}{{\mathcal{N}}_\alpha}           % Nonvan. univ. |a1| = alpha
\newcommand{\Nas}{{\mathcal{N}}_\alpha^*}                 % Symmetrized N_alpha
\newcommand{\eps}{\ensuremath{\varepsilon}}           % alternate epsilon
\renewcommand{\phi}{\ensuremath{\varphi}}             % alternate phi
\newcommand{\nin}{\not \in}                           % Not in
\newcommand{\abs}[1]{\ensuremath{\left| #1 \right|}}  % Absolute value
%\newcommand{\Area}{\rm{Area}}
%\newcommand{\e}{\varepsilon}
%%%%%%%%%%%%%%%%%%%%%%%%%%%%%%%%%%%%%%%%%%%%%%%%%%%%%%%%%%%%%%%%%%%%%%%%%%%%%%%%%%%

\relpenalty=10000
\binoppenalty=10000

\begin{document}

\DOIsuffix{theDOIsuffix}
%%
%% issueinfo for the header line
\Volume{vvv}
\Month{mm}
\Year{2008}

%%    First and last pagenumber of the article. If the option
%%    'autolastpage' is set (default) the second argument may be left empty.
\pagespan{1}{}
%%
%%    Dates will be filled in by the publisher. The 'reviseddate' and
%%    'dateposted' (Published online) entry may be left empty.
\Receiveddate{XXXX}
\Reviseddate{XXXX}
\Accepteddate{XXXX}
\Dateposted{XXXX}
%%
\keywords{omitted area problem, logarithmic capacity, univalent function, symmetrization, local variation}
\subjclass[msc2000]{30C70} 

\title[Iceberg-type Problems]{ Iceberg-type Problems: Estimating Hidden Parts of a Continuum from the Visible Parts}  %

%% Please do not enter footnotes or \inst{}-notes into the optional
%% argument of the author command. The optional argument will go into
%% the header.  If there is only one address the marker \inst{x} may be
%% omitted.

%% Information for the first author.
\author[R.W. Barnard]{Roger W. Barnard}
%% Information for the second author
\author[K. Pearce]{Kent Pearce
  \footnote{Corresponding author:  
  					E-mail: \scriptsize{\textsf{kent.pearce@ttu.edu}},
            Phone: \,806\,742\,2566,
            Fax: \,806\,742\,1112}}
%%
%%    Information for the third author
\author[A.Y. Solynin]{Alexander Yu. Solynin
	\footnote{Partially supported by NSF grant DMS-0525339}}

\address{Department of Mathematics and Statistics,  Texas Tech
University, Box 41042, Lubbock, TX 79409}

\begin{abstract}{We consider the complex plane $\mathbb{C}$ as a space
filled by two different media, separated by the real axis
$\mathbb{R}$. We define $\mathbb{H}_+=\{z: \,\Im \,z>0\}$ to be the upper
half-plane. For a planar body $E$ in $\mathbb{C}$, we discuss a problem of 
estimating characteristics of the ``invisible'' part,
$E_-=E\setminus \mathbb{H}_+$, from characteristics of the whole
body $E$ and its ``visible'' part, $E_+=E\cap \mathbb{H}_+$. In this
paper, we find the maximal draft of $E$ as a function of the
logarithmic capacity of $E$ and the area of $E_+$.
 %AMS No. 30
 }
\end{abstract}


\maketitle


%-----  Section 1:  Introduction  --------------------------------------

\section{Introduction} \label{Introduction}

We will discuss problems (called {\it iceberg-type problems} below) of estimating
characteristics of the ``invisible" part of a compact set $E$ in the complex plane 
$\mathbb{C}$ from some known characteristics of the whole set and its ``visible" part.
We emphasis from the beginning that the problems we study in this paper are not directly
related to (real) physical icebergs.  The problem name reflects the fact that the 
object under consideration consists of two parts, hidden and visible, and the question
is to recover some of the properties of the hidden part from the visible part. During
the ages a titanic work has been done to solve this problem in its
everyday physical setting.

In this paper, we study iceberg-type problems in two-dimensional
space, which will be identified as the complex plane $\mathbb{C}$.
Accordingly, $\overline{\mathbb{C}}=\mathbb{C}\cup\{\infty\}$,
$\mathbb{H}_+=\{z:\,\Im\,z> 0\}$, and $\mathbb{H}_-=\{z:\,\Im\,z<0\}$
will denote the extended complex plane, the upper half-plane, and
the lower half-plane, respectively. The real axis $\mathbb{R}$
will play the role of the surface of interface between
$\mathbb{H}_+$ and $\mathbb{H}_-$.

%\smallskip

For any given compact set $E$ in $\mathbb{C}$, we define $E_+=E\cap
\mathbb{H}_+$ and $E_-=E\cap \overline{\mathbb{H}}_-$. The sets $E_+$ and $E_-$ denote the
visible and hidden parts of $E$, respectively.

%\smallskip

An accumulative characteristic of any body $E$ surrounded by media
is its potential or capacity. In our two-dimensional setting, the
\textit{logarithmic capacity} will be chosen as
the primary characteristic of $E$. We remind the reader that
the logarithmic capacity, $\CAP E$, of a compact set $E$ is defined
by %
$$ %
-\log \CAP E =\lim_{z\to \infty}(g(z)-\log|z|),  %
$$  %
where $g(z)$ denotes Green's function of the unbounded component
$D(E)$ of $\overline{\mathbb{C}}\setminus E$ having singularity at
$z=\infty$. Let ${\mathcal{F}}$ be the collection of all continua
($=$ connected compact
sets) $E$ in $\mathbb{C}$ such that  %
$$ %\be  \label{1.1} %
\CAP(E)=1.  %
$$  %\ee  %

\smallskip

For the measured characteristic of a visible part $E_+$ we will choose
the mass of $E_+$, which, assuming homogeneity of $E$, is
proportional to the area of $E_+$. For $E$ in ${\mathcal{F}}$, the
well known estimates of the
logarithmic capacity show that %
$$ %\be  \label{1.2}  %
0\le \Area(E_+)\le \Area(E)\le \pi (\CAP (E))^2=\pi.  %
$$  %\ee  %



Characteristics of the hidden part $E_-$ which one may want to
control and which are of a particular importance, include: the
draft of the iceberg $H(E)$, the width of the invisible part of
the iceberg $w(E)$, and the safe distance from the iceberg $d(E)$.
Figure~\ref{Figure_Iceberg} illustrates these characteristics
while the precise definitions are as follows:

\be \label{equation draft}  %
H(E)=\max \,(-\Im (z)),   %
\ee %
where the maximum is taken over all $z$ in $E$,  %
\be \label{equation width}  %
w(E)=\max \,(\Re (z_2-z_1)),  %
\ee %
where the maximum is taken
over all $z_1,z_2$ in $E_-$, and %
\be \label{equation safe}  %
d(E)=\max \,(\Re (z_2))-\sup\,(\Re(z_1)),   %
\ee %
where the maximum is taken  over all $z_2$ in $E_-$ and the supremum
is taken over all $z_1$ in $E_+$.

\smallskip

\begin{figure}
$$\includegraphics[scale=.3,angle=0]{iceberg} $$
\caption{Two-dimensional iceberg.}
\label{Figure_Iceberg}
\end{figure}

\smallskip

Then, the extremal problem for each of the functionals
(\ref{equation draft}), (\ref{equation width}), and (\ref{equation safe})
is to find its maximal value over the class ${\mathcal{F}}$
and describe all possible extremal continua.  We define %
\be  \label{equation maximals}  %
{\textbf{(a)}} \ \ \ H({\mathcal{F}})=\max H(E) \quad \quad
{\textbf{(b)}} \ \ \ w({\mathcal{F}})=\max w(E)
\quad \quad {\textbf{(c)}} \ \ \ d({\mathcal{F}})=\max d(E),  %
\ee  %
where in each case the maximum is taken over all sets $E$ in
${\mathcal{F}}$.

Our main goal in this paper is to give a complete solution to
problem (\ref{equation maximals})\textbf{(a)}.  Problems
(\ref{equation maximals})\textbf{(b)} and
(\ref{equation maximals})\textbf{(c)} along
with some other questions will be discussed in the last section.

\medskip


As is well known, problems on the logarithmic capacity of continua
can be reformulated as problems about functions in the class $\Sigma'$ of univalent functions %
\be  \label{equation Sigma} %
f(z)=z^{-1}+a_0+a_1z+\cdots, %
\ee  %
which are analytic in the unit disk $\mathbb{D}$, except for a
simple pole at $z=0$. For $f$ in $\Sigma'$, define
$E_f=\overline{\mathbb{C}}\setminus f(\mathbb{D})$ and define 
$\Sigma'_0=\{f\in \Sigma':\,0\in E_f\}$.

\smallskip

We will solve problem (\ref{equation maximals})\textbf{(a)} by
solving its reformulated dual problem for the class $\Sigma'_0$. 
There is a technical advantage in shifting to the dual problem in that 
the analytical and constructional difficulties which surround the dual 
problem are more tractable than those in the original setting. The 
precise formulation of the dual of the
maximal draft problem (\ref{equation maximals})\textbf{(a)} is the
following problem on the maximal omitted area for the class
$\Sigma'_0$.
For any given real $h$ such that $0<h<4$, find  %
\be  \label{equation dual} %
A(h):=\max \Area (E_f\cap \{w:\,\Re\,w>h\}),  %
\ee  %
where the maximum is taken over all $f$ in $\Sigma'_0$, and find all
functions $f$ in $\Sigma'_0$ extremal for (\ref{equation dual}).
Thus, the question is, for any given $h$, such that $0<h<4$, to maximize the area
omitted by the functions $f$ in $\Sigma'_0$ in the half-plane
$\mathbb{H}_h:=\{w:\,\Re w>h\}$.

We note here that our parameter $h$, which is equal to
the horizontal distance from $w=0$ to the half-plane
$\mathbb{H}_h$, gives as well the value of the maximal draft of
icebergs with visible area $A=A(h)$. In addition, in Corollary~\ref{Corollary
1} (below) we show that the extremal configuration for problem 
(\ref{equation maximals})\textbf{(a)} coincides with the extremal 
configuration for problem (\ref{equation dual}) up to rotation and translation.

For convenience we define $A_f(h) = \Area (E_f\cap \mathbb{H}_h)$.
The maximal omitted area problem (\ref{equation dual})  is solved by the
  following  theorem. %


\bt \label{Theorem 1}  %
Let $h$ satisfy $0<h<4$ and let $f$ belong to $\Sigma'_0$. Then, %
\be \label{equation Thm1-A}  %
A_f(h) \le \pi \beta^2 - 2\beta hr(1-r^2)
\int_{\tau}^1
\left(\frac{t(1-t^2)\sqrt{1-\tau^2t^2}}{(r^2+t^2)^2\sqrt{t^2-\tau^2}}-
\frac{(1-t^2)\sqrt{t^2-\tau^2} }{t(1+r^2t^2)^2\sqrt{1-\tau^2
t^2}}\right)\, dt,
\ee  %
where $r=r(h)$ is the solution to the equation  %
\be \label{equation Thm1-B}  %
h = 2\beta r(1-r^2) \int_0^{\tau} \frac{t(1-t^2)\sqrt{1-\tau^2
t^2} }{(r^2+t^2)^2\sqrt{\tau^2-t^2} }\, dt,
\ee  %
which is unique for $0 < r < 1$ and where %
\be \label{equation Thm1-C}  %
\tau = \sqrt{\frac{(1+r^2)\sqrt{2 - 2r^2 + r^4} - (1+r^4)}{1+3r^2}}  %
\ee  %
and %
\be  \label{equation Thm1-D} %
\beta = \frac{4r\sqrt{r^2 + \tau^2)}}{ (1 + r^2)^2  \sqrt{1 +
\tau^2 r^2)} }.  %
\ee  %

Equality occurs in (\ref{equation Thm1-A}) if and only if
$f=f_{h}$ with $f_h(z)=F(\psi_r^{-1}(z))$ with $r$ defined by (\ref{equation Thm1-B}), where %
\be  \label{equation Thm1-E}  %
z = \psi_r(s) = \frac{ (1-r^2)s - r(1-s^2)}{(1-r^2)s + r(1-s^2)},
\quad \quad s\in \mathbb{D}_+:=\{s\in \mathbb{D}:\, \Re\,s>0\},  %
\ee  %
maps the semidisk  \, ${\mathbb{D}}_+$ conformally onto the unit
disk
$\mathbb D$ and %
\be   \label{equation Thm1-F}   %
F(s) = -2\beta r(1-r^2)\int_0^s \frac{t(t^2+1)\sqrt{1+\tau^2
t^2}}{(t^2-r^2)^2\sqrt{t^2+\tau^2}}\,dt  %
\ee  %
 with the principal branches of the radicals and
 with $\tau$ and $\beta$ defined by (\ref{equation Thm1-C}) and (\ref{equation Thm1-D}).
\et  %

\smallskip

Theorem~\ref{Theorem 1} shows that the maximal omitted area $A(h)$
is given by the explicit expression in the right-hand side of
(\ref{equation Thm1-A}) with $r$, $\tau$, and $\beta$ defined by
(\ref{equation Thm1-B}), (\ref{equation Thm1-C}), and (\ref{equation Thm1-D}),
respectively. Its graph shown in Figure~\ref{Figure_Extremal_Area}
suggests, and we will prove this in Lemma~\ref{Lemma 2.1} in
Section~\ref{Section 2}, that $A(h)$ strictly decreases from $\pi$
to $0$ as $h$ runs from $0$ to $4$. Therefore, the inverse,
$h=\Psi(A)$, of the function $A(h)$ is well defined on $0\le A\le
\pi$.

\bc \label{Corollary 1} %
If $E\in {\mathcal{F}}$ has the visible area $A$, $0<A<\pi$, i.e.
if $\Area (E_+)=A$, then the draft of $E$ is restricted by  %
\be \label{equation Cor1}   %
H(E)\le \Psi(A), %
\ee  %
where the function $\Psi$ is defined above.

Equality occurs in (\ref{equation Cor1}) if and only if $E$
coincides with the continuum $i(E_{f_h}-h)$, where $h=\Psi(A)$ 
and $f_h$ is given in Theorem~\ref{Theorem 1}, up to a horizontal drift.  %
\ec%

\smallskip


The extremal shapes $E(h)=E_{f_h}$ for some typical values of $h$
are displayed in Figure~\ref{Figure_Extremal_Shapes}. As in
previous works on this subject (see \cite{AShS1}-\cite{BS}) the
boundary $\partial E(h)$ consists of the so-called {\it free
boundary} $L_{fr}$ that is an open Jordan arc in $\mathbb{H}_h$
having its ends at the points $h\pm ia$ for some $a$, where $0<a<4$ and the
{\it non-free boundary} $L_{nf}$, which consists of a horizontal
segment $[0,h]$ and two vertical segments $[h,h+ia]$ and
$[h,h-ia]$, see Figure~\ref{Figure_Extremal_Shapes}. The precise
definitions will be postponed until Section~\ref{Section 2}.




\begin{figure}
$$\includegraphics[scale=.4,angle=-90]{graph-of-area-of-h-alt} $$
\caption{Maximal omitted area $A(h)$.}
\label{Figure_Extremal_Area}
\end{figure}

\smallskip


The function $z=\psi(s)$ defined by (\ref{equation Thm1-E}) with
$0<r<1$ maps the semidisk $\mathbb{D}_+$ conformally onto
$\mathbb{D}$ such that $\psi(r)=0$. This reveals the role of the
parameter $r$. The parameters $\tau$ and $\beta$ defined by
(\ref{equation Thm1-C}) and (\ref{equation Thm1-D}) also have
special meanings. Namely, the function $F(s)$ maps the segments
$[-i,-i\tau]$ and $[i\tau,i]$ onto the vertical segments of the
boundary of the corresponding extremal configuration, see
Figure~\ref{Figure_Extremal_Shapes}. In addition, we will prove in
Section~\ref{Section 2} for an extremal function $f_h$ that
$|f'_h(e^{i\theta})|=\beta$ for all $e^{i\theta}$ in the free arc
$l_{fr}=f_h^{-1}(L_{fr})$.

\smallskip

To prove Theorem~\ref{Theorem 1}, we apply techniques developed in
\cite{BS}, \cite{BPS1}, and \cite{BRS1}. These techniques use
symmetrization-type transformations to prove some a priori
smoothness of the boundary, which in turn allows us to apply
Julia-type local variations to find boundary values of the
extremal function. To show that the extremal function can be
recovered from its boundary values and is unique for every $h$, we
prove in Section~\ref{Section 4} several monotonicity lemmas. In
each case, we use a Sturm sequence argument as an essential tool
in our proofs.

We want to mention two other alternative methods, which may work
in the omitted area problems studied in this paper. The first
method is based on the Alt-Caffarelli variational technique which
was developed by J.~Lewis in \cite{L}. His approach does not
require any a priori smoothness and has been found to be very
efficient for omitted area problems, see \cite{L} and \cite{BL}.
The second approach, which was applied in \cite{AShS2} and
\cite{BRS2}, uses Steiner symmetrization to reduce the problem to
the class of typically-real functions. Then, the well-known
integral representation for this class could be used to characterize
the extremal functions.


\section{Extremal configurations and functions} \label{Section 2}%
\setcounter{equation}{0}%

In this section, we collect preliminary results about existence
and geometric properties of extremal functions and configurations.  For notational
convenience, we define $\DD_r(w) = \{z : |z-w|<r \}$, with $\DD_r = \DD_r(0)$, and 
$l_x = \{z : \Re \,z=x \}$; that is, $\DD_r(w)$ is the disk centered at $w$ of radius $r$
and $l_x$ is the vertical line through the point $x$ on the real axis.

\bl \label{Lemma 2.1}  %
\textbf{(a)} For every $h$, where $0\le h\le 4$, there exists $f$ in
$\Sigma'_0$ such that $A_f(h)=A(h)$. In addition, $A(h)$ is
continuous and strictly decreasing in  $0\le h\le 4$.

\textbf{(b)} If $f$ is extremal for $A(h)$, then
$E_f=\mathbb{C}\setminus f(\mathbb{D})$ possesses Steiner symmetry
with respect to $\mathbb{R}$ and circular symmetry with respect to the ray
$\mathbb{R}_0 := \{z : \Re\,z \ge 0 \}$.

\textbf{(c)} For $0<h<4$, the boundary $\partial E_f$ consists of
a free boundary $L_{fr}$ and non-free boundary $L_{nf}$. The
non-free boundary $L_{nf}$ consists of a horizontal segment
$I(h)=[0,h]$ and two vertical segments (possibly degenerate)
$v_f^+=[h,h+ia_f]$ and $v_f^-=[h,h-ia_f]$ with some $0\le a_f<4$
depending on $f$.

The free boundary $L_{fr}$ is an open Jordan rectifiable arc
in $\mathbb{H}_h$ joining the points $h\pm ia_f$.  %
In addition, $\hat L=L_{fr}\cup [-ia_f,ia_f]$ is a closed Jordan
curve that satisfies the following Lavrent'ev condition:
\be \label{equation Lem1} %
{\rm{length}}(J(w_1,w_2))\le C|w_1-w_2| \quad \quad {\mbox{for
$w_1,w_2$ in $\hat L$,}} %
\ee  %
where $C$ is a constant independent of $w_1,w_2$ and $J(w_1,w_2)$
denotes the shortest arc of $\hat L$ between $w_1$ and $w_2$.
\el  %


{\it Proof.}  {\textbf{(a)} Since the omitted area functional
$A_f(h)$ is upper semi-continuous, the existence of an extremal
function, at least one for each $h$, follows from the compactness
of the class ${\Sigma'}_0$. Since $E_f\subset \{w:\,|w|\le 4\}$ for
all $f$ in $\Sigma'_0$, a similar compactness argument easily implies
the continuity of $A(h)$.

Since for any given $f$ in $\Sigma'_0$, the area $A_f(h)$ does not
increase in $0\le h\le 4$, the non-strict monotonicity of $A(h)$
is obvious. Let $0\le h_1<h_2\le 4$. Then, it follows from the
property of the free boundary in part \textbf{(c)}, which is
proved below, that if $f$ is extremal for $A(h_2)$, then $f$ can
not be extremal for $A(h_1)$. Therefore, $A(h)$ is strictly
decreasing on $0\le h\le 4$.

\textbf{(b)}  Symmetry properties can be established via a
standard argument using appropriate Steiner and circular
symmetrizations, cf. \cite{BS}, \cite{BPS1}, \cite{BRS1}.

\textbf{(c)}  Symmetry properties of the extremal configurations
together with the subordination principle, see \cite{H}, imply the
assertion about the non-free boundary $L_{nf}$.

To rule out the case that $L_{fr}$ consists of multiple arcs in
$\mathbb{H}_h$ having their ends on the real axis, we apply
polarization. For the definition and properties of this
transformation the reader may consult \cite{D}, \cite{S},
\cite{BS}.

Let $f$ in $\Sigma'_0$ be an extremal for $A(h)$ and let
$p=\max_{w\in E_f} \,\{\Re \,w\}$ and $p_h=(p+h)/2$. For real
$\tau$, let $E_{f,\tau}^+=E_f\cap \overline{\mathbb{H}}_\tau$,
$E_{f,\tau}^-=E_f\setminus \mathbb{H}_\tau$, and let
$E_{f,\tau}^*$ denote the set symmetric to $E_{f,\tau}^+$ w.r.t.
the vertical line $l_\tau=\{w:\,\Re \,w=\tau\}$. We claim that $E_f$
satisfies the following polarization property (cf. \cite{BS}): %
\be \label{equation polar} %
E_{f,\tau}^*\subset E_{f,\tau}^- \quad \quad {\mbox{for all
$\tau$ such that $p_h\le \tau<p$.  }} %
\ee  %

Indeed, if $E^*_{f,\tau}\not\subset E_{f,\tau}^-$ for some $\tau$,
$p_h\le \tau<h$, then $\widehat{E}_{f,\tau}^p\not=E_f$, where
$\widehat{E}_{f,\tau}^p$ denotes the polarization of $E_f$ into
the half-plane $\mathbb{H}_h^-$. It is also obvious that
$\widehat{E}_{f,\tau}^p$ is not a reflection of $E_f$ in the line
$l_\tau$. Then, the principle of polarization implies the following strict inequality %
$$ %\be \label{equation 2.3}  %
\CAP\,(\widehat{E}_{f,\tau}^p)<\CAP E_f, %
$$  %\ee  %
which easily leads to a contradiction to our assumption that $f$
is extremal for $A(h)$.

Now the symmetry properties of $E_f$ when combined with the
polarization property (\ref{equation polar}) show that $L_{fr}$ is
an open Jordan arc in $\mathbb{H}_h$ joining the points $h\pm
ia_f$ with some $0\le a_f<4$ depending on $f$.

Next, since $E_f$ is bounded, Steiner symmetric w.r.t.
$\mathbb{R}$, and circularly symmetric w.r.t. $\mathbb{R}_0$, the
proof of Lemma~2.2 in \cite{BRS1} shows that $\hat L$ satisfies
the Lavrent'ev condition (\ref{equation Lem1}). In particular,
$\hat L$ and therefore $L_{fr}$ are rectifiable.  The proof of the
lemma is complete. \hfill $\Box$

\medskip


\begin{figure}
$$
\includegraphics[scale=.375,angle=-90]{iceberg-h-0p4-alt}
\includegraphics[scale=.375,angle=-90]{iceberg-h-1p5-alt}
\includegraphics[scale=.375,angle=-90]{iceberg-h-3p0-alt}
$$
\caption{Extremal shapes $E(h)$.}
\label{Figure_Extremal_Shapes}
\end{figure}

\medskip

For $f$ in $\Sigma'_0$ which is extremal for $A(h)$, we define
$l_{fr}=\{e^{i\theta}:\,|\theta|<\theta_1\}$ be the ``free arc'';
that is, $l_{fr}$ is the preimage of $L_{fr}$ under the mapping
$f$. Similarly, we define $l_v^\pm =f^{-1}(v_f^\pm)$ and $l_h^\pm
=f^{-1}(I^\pm(h))$, where $I^+(h)$ and $I^-(h)$ denote,
respectively, the upper part and the lower part of the segment
$I(h)$. We also define $e^{\pm i\theta_1}=f^{-1}(h\pm ia_f)$ and $e^{\pm
i\theta_2}=f^{-1}(h\pm i0)$. Also, we define $l_v = l_v^+ \cup l_v^-$
and $l_h = l_h^+ \cup l_h^-$.

\bl \label{Lemma 2.2} %
For a given $h$, $0<h<4$, let $f$ in $\Sigma'_0$ be extremal for
$A(h)$. Then, there exists a positive $\beta$ such that 

\textbf{(a)} For every sufficiently small positive $\varepsilon$, $f'$ is
bounded on the compact set
$\overline{\mathbb{D}}\setminus\{\mathbb{D}_{\varepsilon}\cup
\mathbb{D}_\varepsilon(e^{i\theta_2})\cup
\mathbb{D}_\varepsilon(e^{-i\theta_2})\}$;

\textbf{(b)} $|f'(z)|=\beta$ if $z\in l_{fr}$;

\textbf{(c)} $|f'(e^{i\theta})|$ strictly increases from $\beta$
to $\infty$ as $\theta$ runs from $\theta_1$ to $\theta_2$;

\textbf{(d)} The vertical non-free boundary is not degenerate,
i.e. $a_f>0$;

\textbf{(e)} $|f'(z)|\to \beta$ as
$z\to e^{i\theta_1}$ such that $z\in \overline{\mathbb{D}}$.  %
\el  %

\textit{Proof.} %
\textbf{(a)} First we prove that $f'$ is bounded near $l_{fr}$. If
not then there is $e^{i\theta_0}$ in $l_{fr}$ and a sequence $z_k\to
e^{i\theta_0}$ such that $z_k\in \mathbb{D}$ for all $k$ in
$\mathbb{N}$ and $f'(z_k)\to \infty$.

Let $\varphi_k$ denote the conformal mapping from $\mathbb{D}$
onto the domain $\mathbb{D}\setminus
\overline{\mathbb{D}}_{\varepsilon_k}(z_k)$ with
$\varepsilon_k=1-|z_k|$ normalized by $\varphi_k(0)=0$,
$\varphi'_k(0)>0$ and define $f_k=\beta_k f \circ \varphi_k$ with
$\beta_k=1-\pi^2\varepsilon^2_k/6$. One can easily verify (see,
for example, Lemma~3.1 in \cite{BS}) that $f_k\in \Sigma'_0$.
Since $0<h<4$ and  $\diam (E_{f_k})\le 4$ by the well-known
Faber's
inequality, an elementary geometric estimate gives %
\be \label{equation Faber} %
\Area (E_{f_k}\cap \mathbb{H}_h^-)\le 4\pi^2\varepsilon_k^2.  %
\ee  %

Using (\ref{equation Faber}) and the mean value property of the
subharmonic function
$|f'(z)|^2$, we can estimate the area $A_{f_k}(h)$ as follows:  %
\begin{alignat}{10}  \label{equation AreaA} %
A_{f_k}(h)&=&\beta_k^2\left(A_f(h)+ \Area
(f(\mathbb{D}_{z_k}(\varepsilon_k)))\right)-\Area(E_{f_k}\cap\mathbb{H}_h^-)\
\ \ \ \ \ \ \ \ \ \ \ \ \ \ \ \ \ \ \ \ \ \ \ \ \ \ \ \ \
\\
&\ge& \beta_k^2\left(A_f(h)+\pi \varepsilon^2_k|f'(z_k)|^2\right)-
4\pi^2\varepsilon_k^2\ge
A_f(h)+\left(\pi|f'(z_k)|^2-C\right)\varepsilon_k^2+o(\varepsilon_k^2),\nonumber%
\end{alignat}  %
with some constant $C>0$ independent of $f$. Since $f'(z_k)\to
\infty$ as $k\to \infty$, (\ref{equation AreaA}) contradicts the
extremality of $f$. Therefore, $f'$ is bounded near $l_{fr}$.

\smallskip


\textbf{(b)} First we show that $|f'(z)|$ is constant a.e.\ on
$l_{fr}$. Since $L_{fr}$ is Jordan locally rectifiable, it follows
that the
non-zero finite limit %
\be   \label{equation Lem2-b-1}   %
f'(\zeta)=\lim_{z\to\zeta,z\in \overline{\mathbb{D}}}
\frac{f(z)-f(\zeta)}{z-\zeta}\not= 0,\infty %
\ee   %
exists  a.e. on $l_{fr}$; see \cite[Theorem 6.8, Exercise
6.4.5]{P}. Assume that %
\be   \label{equation Lem2-b-2}%
0<\beta_1=|f'(e^{i\nu_1})|<|f'(e^{i\nu_2})|=\beta_2<\infty  %
\ee     %
for $e^{i\nu_1}, e^{i\nu_2}\in l_{fr}$. Note that (\ref{equation
Lem2-b-1}) and (\ref{equation Lem2-b-2}), combined with the fact
that $f'$ is bounded near $l_{fr}$, allow us to apply the
two-point variational formulas, see \cite[Lemma~10]{BS} or
\cite[Lemma~5]{BPS1}. Namely, for fixed positive $k_1$, $k_2$ such
that $ 0<k_1<1<k_2$ and $k_1\beta_1^{-1}>k_2\beta_2^{-1}$ and
fixed $\varphi>0$ small enough, we consider the two-point
variation $\tilde D$ of $D=f(\mathbb{D})$ centered at
$w_1=f(e^{i\nu_1})$ and $w_2=f(e^{i\nu_2})$ with inclinations
$\varphi$ and radii $\e_1=k_1 \e$, $\e_2=k_2 \e$, respectively;
see \cite[Section~2]{BPS1}. Computing the change in the area by 
\cite[formula (2.11)]{BPS1}, we find %
\be \label{equation Lem2-b-3}%
{\rm{Area}}\,(\mathbb{C}\setminus\tilde D)-
{\rm{Area}}\,E_f=\frac{2\pi\varphi-\sin 2\pi\varphi}{2\sin^2
\pi\varphi}\e^2(k_1^2-k_2^2)+o(\e^2)>0 %
\ee   %
for all $\e>0$ small enough. Similarly, applying  \cite[formula
(2.10)]{BPS1}, we get  %
\be \label{equation Lem2-b-4}  %
\frac{\CAP(\mathbb{C}\setminus \tilde
D)}{\CAP(E_f)}=\left[\frac{\varphi(2+\varphi)}{6(1+
\varphi)^2}\frac{k_1^2}{\beta_1^2}-
\frac{\varphi(2-\varphi)}{6(1-\varphi)^2}\frac{k_2^2}{\beta_2^2}\right]
\e^2+o(\e^2)<0  %
\ee   %
for all $\e>0$ small enough and $\varphi$ chosen such that the
expression in the brackets is positive.

Inequalities (\ref{equation Lem2-b-3}) and (\ref{equation Lem2-b-4}) lead to
a contradiction to the extremality of $f$ for $A(h)$, via a
standard subordination argument.
  Thus $|f'(e^{i\theta})|=\beta$ a.e.\  on $l_{fr}$ with some $\beta>0$.

To prove that $|f'(e^{i\theta})|=\beta$ everywhere on $l_{fr}$, we consider the
auxiliary conformal mapping %
\be \label{equation 2.7}  %
g=\varphi\circ f\circ k_\tau \quad \quad {\mbox{ with
\ \ $\varphi(w)=1/(w-p_h)$,}}  %
\ee  %
 where $p_h$ is defined in the proof of
Lemma~\ref{Lemma 2.1}, and with %
$$  %\be \label{equation 2.4}  %
k_\tau(\zeta)=k^{-1}(\tau k(\zeta)), \quad {\mbox{ where
$k(\zeta)=\zeta/(1-\zeta)^2$ and $\tau=1/\sin^{2}(\theta_2/2)$.}}  %
$$  %

We note that $k_\tau$ maps the slit disk
$\mathbb{D}'=\mathbb{D}\setminus[-1,-r_0]$, where
$r_0=(\sqrt{\tau}-\sqrt{\tau-1})^2$, conformally and one-to-one
onto $\mathbb{D}$ in such a way  that the radial slit is mapped
onto the arc $l_h=\{e^{i\theta}:\,|\theta-\pi|\le \pi-\theta_2\}$.


Let $D_g'=g(\mathbb{D}')$ and let $D_g=D'_g\cup
((p_h-h)^{-1},-p_h^{-1}]$. By the Schwarz reflection principle,
the function $g$ can be continued to a function, still denoted by
$g$, which maps the whole disk $\mathbb{D}$ conformally and
one-to-one onto $D_g$. It follows from Lemma~\ref{Lemma
2.1}\textbf{(c)} that $D_g$ is a bounded Jordan domain, whose
boundary satisfies the Lavrent'ev condition (\ref{equation Lem1})
for some $C>0$. Therefore, $D_g$ is a Smirnov domain; see
\cite[Sections~7.3, 7.4]{P}.
Thus, $\log |g'|$ can be represented by the Poisson integral %
\be  \label{equation Poisson}%
\log |\varphi'(w)f'(z)k'_\tau(\zeta)|=\log|g'(\zeta)|=\frac{1}{2\pi}\int_0^{2\pi}P(r,\psi-t)%
\log|g'(e^{it})|\, dt %
\ee   %
with boundary values defined a.e. on $\T$; see \cite[p. 155]{P}.
Equation (\ref{equation Poisson}) easily implies that
$$|g'(e^{i\psi})|=\beta |\varphi'(f(k_\tau(e^{i\psi}))||k'_\tau(e^{i\psi})|$$ 
for all $e^{i\psi}$ such that $k_\tau(e^{i\psi})\in l_{fr}$ and therefore
$|f'(e^{i\theta})|=\beta$ for all $e^{i\theta}\in l_{fr}$.  In
addition, (\ref{equation Poisson}) implies that $\log f'$ is
bounded on $\overline{\mathbb{D}}$ outside any neighborhoods of
the points $z=0$, $z=-1$, and $z=e^{\pm i\theta_2}$.

\smallskip



\textbf{(c)} Since $E_f$ is Steiner symmetric w.r.t.\  $\R$, the
strict monotonicity of $|f'|$ along $l_v^+$ follows from
\cite[Lemma~4]{BS}. To prove that $|f'(e^{i\theta})|>\beta$ for
all $e^{i\theta}\in l_v\setminus\{0\}$, we assume that
$\beta=|f'(e^{i\nu_1})|>|f'(e^{i\nu_2})|=\beta_2$ with
$e^{i\nu_1}\in l_{fr}$ and some $e^{i\nu_2}\in l_v^+$.  Then,
applying the two-point variation as above, we get inequalities
(\ref{equation Lem2-b-3}) and (\ref{equation Lem2-b-4}), contradicting the
extremality of $f$ for $A(h)$, again via a subordination argument.
Hence, $|f'(e^{i\theta})|\ge \beta$ for all $e^{i\theta}\in l_v$
which, when combined with the strict monotonicity property of
$|f'|$, leads to the strict inequality $|f'(e^{i\theta})|>\beta$
for $e^{i\theta}\in l_v$.

\smallskip

\textbf{(d)} Assume that $a_f=0$. Then, $\theta_1=\theta_2$,
$L_{nf}=I(h)$, and $\hat L=L_{fr}\cup \{h\}$. In addition,
$|f'(e^{i\theta})|=\beta>0$ for all $e^{i\theta}\in l_{fr}$ by
part~\textbf{(b)} of this proof.

In the notation of part \textbf{(b)}, we consider the function
$g=\varphi\circ f\circ k_\tau$ defined by (\ref{equation 2.7}),
which maps $\mathbb{D}$ conformally onto the domain $D_g$. As we
have mentioned above,
 $\log |g'(\zeta)|$ can be represented
by the Poisson integral (\ref{equation Poisson}).%

Since $|k'_\tau(e^{i\psi})|\to 0$ as $\psi\to \pi$, it follows
that
$|g'(e^{i\psi})|=\beta|\varphi'(k_\tau(e^{i\psi}))||k'_\tau(e^{i\psi})|\to
0$ as $\psi\to \pi$. Therefore, %
\be \label{equation 2.5} %
\log|g'(e^{i\psi})|\to -\infty \quad \quad {\mbox{as $\psi\to
\pi$.}}  %
\ee  %

From  (\ref{equation Poisson}) and (\ref{equation 2.5}), using the
well-known properties of the radial limits of the Poisson
integral, we  obtain that %
\be \label{equation 2.6} %
\log|g'(-r)|\to -\infty \quad \quad {\mbox{as $r\to 1^-$.}}  %
\ee  %



Now we show that $g$ has a finite non-zero angular derivative at
$\zeta=-1$. To do this, we construct two comparison functions
$f_1$ and $f_2$. Let $f_1$ map $\mathbb{D}$ conformally onto the
vertical strip $\{w:\,0<\Re\,w<h\}$ such that $f_1(0)=h/2$,
$f_1(-1)=h$ and let $g_1=\varphi\circ f_1$. Then, of course,
$g_1'(-1)$ exists and $g_1'(-1)\not=0,\infty$. Since $g_1(\mathbb{D})\subset
g(\mathbb{D})$ and $g_1(-1)=g(-1)=1/(h-p_h)$, we can apply the
comparison Theorem~4.14 in \cite{P} to conclude that $g$ has the
angular
derivative $g'(-1)$ and %
\be  \label{equation Lem2-e-2} %
|g'(-1)|=c_1|g'_1(-1)| \quad {\mbox{where $0\le c_1<\infty$.}}  %
\ee  %

Next we construct our second comparison function.  We define
$K_{p_h}=E_f\cap \overline{\mathbb{H}}_{p_h}$ and $K^*_{p_h}$
be the set symmetric to $K_{p_h}$ w.r.t. the vertical line
$l_{p_h}$. Define $\Omega=\mathbb{C}\setminus\left(K_{p_h}\cup
K^*_{p_h}\right)$, let $f_2$ map $\mathbb{D}$ conformally onto
$\Omega$ such that $f_2(0)=\infty$, $f_2(-1)=h$, and let
$g_2=\varphi\circ f_2$. Since the boundary $\partial \Omega$ is
analytic in a vicinity of $w=h$, it follows that
$g'_2(-1)$ exists and $g'_2(-1)\not=0,\infty$.

It follows from equation (\ref{equation polar}) in the proof of
Lemma~\ref{Lemma 2.1}\textbf{(c)} that $g(\mathbb{D})\subset
g_2(\mathbb{D})$. Now, Theorem~4.14 in
\cite{P} implies that %
$$  %
|g'_2(-1)|=c_2|g'(-1)| \quad {\mbox{where $0\le c_2<\infty$.}}  %
$$  %


This together with (\ref{equation Lem2-e-2}) shows that the finite
non-zero angular derivative  $g'(-1)$ exists. Now Proposition~4.7
\cite{P} implies that $g'(\zeta)$ has the finite
angular limit $g'(-1)$ at $\zeta=-1$ where $g'(-1)\not 0$. In particular, %
$$ %
|g'(-r)|\to |g'(-1)|\not=0 \quad \quad {\mbox{as $r\to 1^-$}} %
$$  %
contradicting (\ref{equation 2.6}). This proves that
$a_f>0$.


\textbf{(e)} To show that $|f'|$ is continuous at $e^{\pm
i\theta_1}$, we again use the function $g$ defined by
(\ref{equation 2.7}). Using Theorem~4.14 in \cite{P} with $g_1$
defined in part \textbf{(d)} of this proof as a comparison
function, we conclude that the finite angular derivative
$g'(k_\tau^{-1}(e^{i\theta_1}))$, and therefore the angular
derivative $f'(e^{i\theta_1})$, exists finitely.

By the reflection principle, $f$ can be continued analytically
across $l_v^-$. By Lemma~\ref{Lemma 2.1}, $E_f$ is Steiner
symmetric w.r.t. $\mathbb{R}$ and circularly symmetric w.r.t.
$\mathbb{R}_0$. Using these facts it is not difficult to see that
this analytic continuation, say $\tilde f$, of $f$ is univalent in
the disk
$U=\{z:\,|z-\varepsilon_0e^{i(\theta_1+\theta_2)/2)}|<\rho_0\}$
for a sufficiently small positive $\varepsilon_0$ and
$\rho_0=|e^{i\theta_1/2}-\varepsilon_0e^{i\theta_2/2}|$. By
Proposition~4.9 \cite{P}, the function $\tilde f$ has the angular
derivative ${\tilde f}'(e^{i\theta_1})$ at $z=e^{i\theta_1}$,
which of course coincides with the angular derivative
$f'(e^{i\theta_1})$.

We have $l_v^-\subset U$. Since $|f'(e^{i\theta})|$ is monotone
and greater than $\beta$ on $l_v^-$, it follows that
$\lim_{\theta\to \theta_1^+}
f'(e^{i\theta})=\beta_0e^{-i\theta_1}$ where
$0 < \beta \le \beta_0$. Therefore,  %
\be  \label{equation 2.8} %
f'(z)\to \beta_0e^{-i\theta_1} \quad \quad {\mbox{as \ \ $z\to
e^{i\theta_1}$}}  %
\ee  %
in any Stolz angle in $\mathbb{D}$ with the vertex at
$e^{i\theta_1}$. To show that $\beta_0=\beta$, we use the Poisson
integral (\ref{equation Poisson}). Let
$\psi_1=\arg(k_\tau^{-1}(e^{i\theta_1}))$. If $\beta_0\not=\beta$,
then the theorem about radial limits of the Poisson integral
implies
that %
$$  %
\lim_{r\to 1^-}\log|g'(re^{i\psi_1})|=\frac12\lim_{\varepsilon \to
0}
\log|g'(e^{i(\psi_1+\varepsilon)})g'(e^{i(\psi_1-\varepsilon)})|. %
$$  %
This implies that $|f'(k_\tau(re^{i\theta_1}))|\to
\sqrt{\beta\beta_0}$ as $r\to 1^-$, which together with
(\ref{equation 2.8}) shows that we must have $\beta_0=\beta$.

Using the Poisson integral (\ref{equation Poisson}) once more, we
conclude that $\log|g'(\zeta)|$ is continuous for $\zeta$ such
that $|\zeta|\le 1$ and $|\zeta-k_\tau^{-1}(e^{i\theta_1})|$ is
small enough. Since $g=\varphi\circ f\circ k_\tau$ and $\varphi$
and $k_\tau$ are conformal in the corresponding domains the latter
implies \textbf{(e)}.

 The proof of Lemma~\ref{Lemma 2.2} is complete. \hfill  $\Box$






\section{Closed form of the extremal functions and the proof of Theorem~\ref{Theorem 1}}
\label{Section 3} %
\setcounter{equation}{0}  %

Lemmas \ref{Lemma 2.1} and \ref{Lemma 2.2} provide sufficient
information to find a closed form of the function $f$ extremal for
$A(h)$ when $0<h<4$. It is convenient to work in the
auxiliary $s$-plane with $z=\psi_r(s)$ defined by (\ref{equation Thm1-E}).
We note that this auxiliary mapping was already used in
\cite{AShS2} to solve the minimal area $a_2$-problem for convex
functions.

The function $z=\psi_r(s)$ maps the semidisk ${\mathbb{D}}_+$
conformally onto
$\mathbb{D}$ such that %
$$  %\be  \label{equation 3.1}  %
\psi_r(r)=0,\quad \psi_r(i)=e^{i\theta(r)}, %
$$  %\ee  %
where %
\be  \label{equation arcsin}  %
\theta(r)=2\arcsin \frac{2r}{1+r^2}.  %
\ee  %

\bl \label{Lemma 3.1}  %
Let $f$ be extremal for $A(h)$, $0<h<4$, and let
$F_r(s)=f(\psi_r(s))$, $0<r<1$. Then, there are parameters $r,\tau,\beta$ where $0<r<1$, $0<\tau<1$, and $\beta>0$ such that  %
\be  \label{equation Lem3} %
F'_r(s)=-\frac{2\beta
r(1-r^2)s(s^2+1)(1+\tau^2s^2)^{1/2}}{(s^2-r^2)^2(s^2+\tau^2)^{1/2}} %
\ee  %
with the principal branches of the radicals.  %
\el  %

\textit{Proof.} Let $\theta_1$ and $\theta_2$ be the angles
defined for $f$ as in Section~\ref{Section 2}. Since $\theta(r)$
defined by (\ref{equation arcsin}) strictly increases in $0<r<1$, its
inverse, $r(\theta)$, is well defined. Choose $r=r(\theta_1)$. For
this $r$, let $i\tau=\psi_r^{-1}(e^{i\theta_2})$. Then, $0<\tau<1$.
By Lemma~\ref{Lemma 2.2}, there is a positive $\beta$ such that
$|f'(e^{i\theta})|=\beta$ for all $|\theta|\le \theta_1$.

Let $\Phi(s)=\Phi(s;r,\tau,\beta)$ denote the expression in the
right-hand side of (\ref{equation Lem3}) considered as a function
of $s\in \mathbb{D}_+$ for the values of $r$, $\tau$, and $\beta$
chosen above.

It follows from (\ref{equation Sigma}) and (\ref{equation Thm1-E})
that the limit %
\be  \label{equation limit}  %
\lim_{s\to r}\left(F'_r(s)/\Phi(s)\right) =\frac{4r(r^2+\tau^2)^{1/2}}{\beta(1+r^2)^2(1+r^2\tau^2)^{1/2}} %
\ee  %
exists and is finite and non-zero.  Using (\ref{equation limit}) one
can easily show that the function  %
$$  %\be  \label{equation 3.6}  %
g(s)=u(s)+iv(s):=\log\left(F'_r(s)/\Phi(s)\right)  %
$$  %\ee  %
is analytic and single-valued on $\mathbb{D}_+$.

It follows from Lemma~\ref{Lemma 2.1} and the definition of $\Phi(s)$
that $g(s)$ takes real values on the vertical diameter $[-i,i]$
except its three singularities at the points $s=0$, $s=i\tau$, and
$s=-i\tau$. By the Schwarz reflection principle, $g(s)$ can be
continued as an analytic multi-valued function in the punctured
disk $\mathbb{D}'=\mathbb{D}\setminus\{0,\pm i\tau\}$.

To analyze the nature of the multi-valuedness of $g$, we compute the
periods $\omega_0$, $\omega_1$, and $\omega_{-1}$ of $g$ at the
singularities $s=0$, $s=i\tau$, and $s=-i\tau$, respectively.

Since $F_r(s)$ maps the segments $[-i\tau,0]$ and $[0,i\tau]$, each
one-to-one onto the horizontal segment $[0,h]$, it follows that
$F_r(s)$ is analytic near $s=0$ and its Taylor expansion at $s=0$
has the form $F_r(s)=Cs^2+\cdots$ where $C\not=0$. Then, we
have  %
$$  %
\frac{F''_r(s)}{F'_r(s)}=\frac1s +{\mbox{non-negative powers of
$s$.}}  %
$$  %
Using this, we easily find that %
$$  %\be  \label{equation 3.7}  %
\omega_0=\int_{|s|=\varepsilon}dg(s)=\int_{|s|=\varepsilon}
\left(\frac{F''_r(s)}{F'_r(s)}-\frac{\Phi'(s)}{\Phi(s)}\right)\,ds=0  %
$$  %\ee  %
for all sufficiently small positive $\varepsilon$.  Similarly, we find %
$$  %\be  \label{equation 3.8}  %
\omega_1=\int_{|s-i\tau|=\varepsilon}dg(s)=\int_{|s-i\tau|=\varepsilon}
\left(\frac{F''_r(s)}{F'_r(s)}-\frac{\Phi'(s)}{\Phi(s)}\right)\,ds=0.  %
$$  %\ee  %
By symmetry, we also have $\omega_{-1}=0$.

Since all periods of $g$ are zero, the function $g(s)$ is analytic
and single-valued on $\mathbb{D}$.

\smallskip

We claim that $u(s):=\Re\,g(s)\equiv 0$ on $\mathbb{D}$. To prove
this, we test the boundary values of $u$. For $s=e^{it}$ with
$|t|< \pi/2$, using Lemma~\ref{Lemma 2.2}\textbf{(b)} and
(\ref{equation Thm1-E}) we compute %
$$  %
|F'_r(s)|=2\beta r(1-r^2)
\frac{|s^2+1|}{|s^2-r^2|^2}=|\Phi(s)|,  %
$$  %
which shows that $u(e^{it})=0$ for $|t|< \pi/2$. By
Lemma~\ref{Lemma 2.2}\textbf{(e)}, $|f'(z)|\to  \beta$ as $z\to
e^{i\theta_1}$. Using the explicit expressions for $\psi_r$ and
$\Phi$, see (\ref{equation Thm1-E}) and (\ref{equation Lem3}), we
easily find that $|F'_r(s)/\Phi(s)|\to 1$ as $s\to i$. Thus, $u$
has  boundary value $0$ at $s=i$. By symmetry, $u(e^{it})=0$
everywhere on $\mathbb{T}$.

Since $u$ is harmonic in $\mathbb{D}$ and continuous on
$\overline{\mathbb{D}}$, the maximum principle implies that
$u(s)\equiv 0$ on $\mathbb{D}$.  Then, of course, $g(s)$ is
constant on $\mathbb{D}$, and this constant has the value $0$
since $\Im\,g(r)=0$ by (\ref{equation limit}). This proves the
lemma. \hfill $\Box$


\smallskip

Using the closed form (\ref{equation Lem3}) combined with some
computational results, the proofs of which are postponed until
Section~\ref{Section 4}, we can prove our main theorem.


\smallskip

\textit{Proof of Theorem~\ref{Theorem 1}.} If $0<h<4$, let $f$
be an extremal function for $A(h)$, which exists by
Lemma~\ref{Lemma 2.1}. Let $F_r(s)=f(\psi_r(s))$ be defined as in
Lemma~\ref{Lemma 3.1}. Then, $F'_r(s)$ has the form (\ref{equation Lem3}).

We claim that there is a unique set of parameters $r=r(h)$,
$\tau=\tau(h)$, and $\beta=\beta(h)$, for which the function
$f_h(z)=F_r(\psi_r^{-1}(z))$, with $F_r(s)$ defined by
(\ref{equation Lem3}), is  in $\Sigma'_0$. Then, of course, $f_h$
will be the unique extremal for $A(h)$.

Expanding (\ref{equation Lem3}) into a Laurent series at $s=r$, we
obtain  %
$$  %\be  \label{equation 3.9}  %
F'_r(s)= \frac{A_{-2}}{(s-r)^2}+\frac{A_{-1}}{s-r}+A_0+{\mbox{positive powers of $(s-r)$,}} %
$$  %\ee    %
where  %
\be  \label{equation Laurent coef} %
A_{-2}=-\frac{\beta(1-r^4)(1+r^2\tau^2)^{1/2}}{2(r^2+\tau^2)^{1/2}} %
\ee  %
and  %
$$  %\be  \label{equation 3.11} %
A_{-1}=-\frac{\beta
r(1-r^2)((1+3r^2)\tau^4+2(1+r^4)\tau^2-(1-r^2))}{2(r^2+\tau^2)^{3/2}(1+r^2\tau^2)^{1/2}}.  %
$$  %\ee  %

Since $F_r(s)$ is a single-valued function in $\mathbb{D}_+$, we
must have $A_{-1}=0$. This gives %
$$   %\be  \label{equation 5.6} %
\tau=\sqrt{\frac{(1+r^2)\sqrt{2-2r^2+r^4}-(1+r^4)}{1+3r^2}}, %
$$  %\ee  %
which is equation (\ref{equation Thm1-C}) of Theorem~\ref{Theorem
1}.

To find $\beta$, we use the normalization $\lim_{z\to
0}(-z^2f'(z))=1$. Then, using (\ref{equation Thm1-E}) and (\ref{equation Laurent coef}), we obtain %
$$  %\be \label{equation 5.7} %
1=\lim_{s\to
r}\left(-\psi_r^2(s)\frac{F'(s)}{\psi'_r(s)}\right)=\lim_{s\to
r}\left(\frac{\beta(s-r)^2(1+r^2)^2(1+\tau^2s^2)^{1/2}}{4r^2(s-r)^2(s^2+\tau^2)^{1/2}}\right).  %
$$  %\ee  %
From this we find %
$$  %
\beta=\frac{4r(r^2+\tau^2)^{1/2}}{(1+r^2)^2(1+r^2\tau^2)^{1/2}},  %
$$   %
which is equation (\ref{equation Thm1-D}) of Theorem~\ref{Theorem
1}.



Next, using conditions  $f(e^{i\theta_2})=F_r(i\tau)=h$, we can
find an equation that links $r$ and $h$: %
$$  %
h=\int_0^{i\tau}F'(s)\,ds=2\beta r(1-r^2)\int_0^\tau
\frac{t(1-t^2)(1-\tau^2t^2)^{1/2}}{(t^2+r^2)^2(\tau^2-t^2)^{1/2}}\,dt,  %
$$  %
which is equation (\ref{equation Thm1-B}) of Theorem~\ref{Theorem 1}.

Let $h(r)$ denote the right-hand side of (\ref{equation Thm1-B})
with $\tau$ and $\beta$ considered as functions of $r$ defined by
(\ref{equation Thm1-C}) and (\ref{equation Thm1-D}). In
Lemma~\ref{Lemma h} in Section~\ref{Section 4}, we will show that
$h(r)$ strictly decreases from $4$ to $0$ as $r$ runs from $0$ to
$1$. Therefore, for every $h$, such that $0<h<4$, (\ref{equation Thm1-B}) has a
unique solution $r=r(h)$ whenever $0<r<1$.

Thus, we have proven that for every $0<h<4$, there is a unique
function $f_h$ in $\Sigma'_0$ extremal for $A(h)$. In addition, we
have shown that the derivative
$F'_r(s)=f'_h(\psi_r(s))\psi'_r(s)$, where $r=r(h)$ is defined by
(\ref{equation Thm1-B}), is given by (\ref{equation Lem3}).
Integrating (\ref{equation Lem3}), we obtain (\ref{equation Thm1-F}).

To complete the proof of Theorem~\ref{Theorem 1}, we have to find
the maximal omitted area $A(h)$. This calculation will be given in
Lemma~\ref{Lemma 4.1} below.  \hfill $\Box$




%%%%%%%%%%%%%%%%%%%%%%%%%%%%%%%%%%%%%%%%%%%%%%%%%%%%%%%%%%%%%%%%%%%%%%%%%%%

\section{Area functional and monotonicity lemmas} \label{Section 4}

\setcounter{equation}{0} %

\bl  \label{Lemma 4.1} For $0<h<4$, the maximal omitted area $A(h)$ is given by %
\be \label{equation A(h)}  %
A(h) = \pi \beta^2 - 2\beta hr(1-r^2) \int_{\tau}^1
\left(\frac{t(1-t^2)\sqrt{1-\tau^2t^2}}{(r^2+t^2)^2\sqrt{t^2-\tau^2}}-
\frac{(1-t^2)\sqrt{t^2-\tau^2} }{t(1+r^2t^2)^2\sqrt{1-\tau^2
t^2}}\right)\, dt
\ee  %
with $r$, $\tau$, and $\beta$ defined in Theorem~\ref{Theorem 1}.
\el %

{\textit{Proof.}} Let $f$ be extremal for $A(h)$ and let $F_r(s)$
with $r=r(h)$ be defined for $f$ as in Lemma~\ref{Lemma 3.1}.
Applying the standard line integral formula for the area, we find  %
$$  %\be  \label{equation a} %
A(h)=\frac12 \Im\,\int_{\partial E_f} \bar w\,dw =\frac12 \Im\,
\int_{h+ai}^{h-ai}\bar w\,dw +\frac12 \Im\, \int_{L_{fr}}\bar
w\,dw
 = -ha + \frac12 \Im\, \int_{L_{fr}}\bar w\,dw ,
$$  %\ee  %
 where %
\be  \label{equation a}  %$$  %
a=\Im f(e^{i\theta_1})=\Im\, \int_{i\tau}^i F'_r(s)\, ds = \beta
r(1-r^2) \int_{\tau}^1 \frac{2t(1-t^2)\sqrt{1-\tau^2 t^2
}}{(t^2+r^2)^2\sqrt{t^2-\tau^2 }} \, dt.  %
\ee  %$$  %
Now, taking the condition $|f'(z)|=\beta$ for $z\in l_{fr}$ into
account, we find the integral over the free boundary:
\begin{eqnarray*}
\frac12 \Im\, \int_{L_{fr}}\bar w\,dw &=& \frac12 \Re
\int_{\theta_1}^{-\theta_1} f(e^{i\theta})e^{-i\theta}
{\overline{f'(e^{i\theta})}}\,d\theta = \frac{\beta^2}{2}
\Re\,\int _{\theta_1}^{-\theta_1}\frac{f(e^{i\theta})e^{i\theta}}{e^{2i\theta}f'(e^{i\theta})}\,d\theta \\
  &=& -\frac{\beta^2}{2} \Im\int_{\T}\frac{f(z)}{z^2 f'(z)}\,dz+
%  2\beta^2 h \int_{\theta_1}^{2\pi-\theta_1}\frac{d\theta}{|f'(e^{i\theta})|} \\
   \frac{\beta^2}{2}\Re\,\int_{\theta_1}^{2\pi-\theta_1} \frac{f(e^{i\theta})}{e^{i\theta}f'(e^{i\theta})}d\theta \\
 &=& -\frac{\beta^2}{2} \Im \,{\rm{Res}}\,\left[ \frac{f(z)}{z^2f'(z)},0 \right] +
 \beta^2 h \int_{\theta_1}^{\theta_2}\frac{d\theta}{|f'(e^{i\theta})|} \\
 &=& \pi \beta^2 + \beta^2 h \int_{\theta_1}^{\theta_2}\frac{d\theta}{|f'(e^{i\theta})|}.
\end{eqnarray*}
To find
$\int_{\theta_1}^{\theta_2}|f'(e^{i\theta})|^{-1}\,d\theta$, we
change variables via $z=\psi_r(s)$ to obtain %
$$  %\be \label{equation 4.4}
\int_{\theta_1}^{\theta_2}\frac{d\theta}{|f'(e^{i\theta})|} =
\int_{\tau}^1 \frac{|\psi'(it)|^2}{|F'(it)|} \,  dt =
\frac{2r(1-r^2)}{\beta}\int_{\tau}^1 \frac{t(1-t^2)\sqrt{1-\tau^2
t^2 }}{(t^2+r^2)^2\sqrt{t^2-\tau^2 }} \, dt.  %
$$  %\ee   %
Combining all of these calculations we obtain (\ref{equation A(h)}). \hfill $\Box$

\smallskip

 After integration, an explicit formulation for the maximal
omitted area $A(h)$ can be expressed as a function of
$r$ that is a complicated combination of polynomials, square
roots, and logarithms. Although explicit, this form does not give
us any computational advantages. In contrast, to prove the
monotonicity of the function $h(r)$ defined by (\ref{equation
Thm1-B}), it is useful to express the integral in (\ref{equation
Thm1-B}) in terms of elementary functions. The graph of $h(r)$ is shown in
Figure~\ref{Figure_h_and_a}.



\begin{figure}
$$
\includegraphics[scale=.35,angle=-90]{graph-of-h-alt}
\includegraphics[scale=.35,angle=-90]{graph-of-a-alt}
$$
\caption{Functions $h(r)$ and $a(r)$.}
\label{Figure_h_and_a}
\end{figure}

\medskip


Changing the variables via $t^2 = \tau^2x$, we can rewrite
(\ref{equation Thm1-B}) as %
$$  %
h=\beta r(1-r^2)\tau \int_0^1
\frac{(1-\tau^2x)(1-\tau^4x)}{(r^2+\tau^2x)^2}
\frac{dx}{\sqrt{(1-x)(1-\tau^4x)}}.
$$  %
Expanding the rational function in the integrand into partial
fractions and then integrating, yields the following explicit
representation for $h=h(r)$ as a function of $r$: %
\be  \label{equation h(r)} %
h(r)=\frac{4(1-r^2)\left(P(1+r^2)-(1+r^4)\right)^{1/2}\left(1+3r^2+r^2(P-1+2r^2)Q\right)}{(1+r^2)^2(1+3r^2)^{1/2}
(P-1+2r^2)^{1/2}\left(r^2(P+1)+1-r^4\right)^{1/2}},  %
\ee  %
 where  %
 \be  \label{equation P and Q} %
 P=(2-2r^2+r^4)^{1/2} \quad {\mbox{and}} \quad Q=\log\frac{P(1+r^2)+r^2(3-r^2)}{(1+r^2)(2-P+r^2)}.  %
 \ee  %

\bl  \label{Lemma h} %
 The function $h=h(r)$  defined by (\ref{equation h(r)}) strictly decreases from $4$ to $0$
 as $r$ runs from $0$ to $1$.%
\el  %

\textit{Proof.} %
 Differentiating (\ref{equation h(r)}), we find  %
$$  %
h'(r) = \frac{-16r(1-r^2)\left((c_0 + c_1P)+ (d_0 +
d_1P)Q\right)}{D},  %
$$  %
where  %
\begin{eqnarray*}
c_0 &=&  -96r^{22}+184r^{20}+144r^{18}-318r^{16}-228r^{14}+220r^{12}+296r^{10}\\
    &+& 868r^8-436r^6+84r^4+64r^2-14,\\
c_1 &=& 96r^{20}-88r^{18}-280r^{16}\\
    &+& 34r^{14}+410r^{12}+334r^{10}-50r^8+398r^6-58r^4-38r^2+10,\\
d_0 &=&  -32r^{24}+72r^{22}+16r^{20}-124r^{18}-11r^{16}+127r^{14}+1703r^{12}\\
    &-& 3889r^{10}+4041r^8-1475r^6-347r^4+361r^2-58,\\
d_1 &=& 32r^{22}-40r^{20}-72r^{18}+56r^{16}+111r^{14}+3r^{12}\\
    &+&1535r^{10}-2293r^8+1125r^6+121r^4-235r^2+41, %
\end{eqnarray*}  %
and %
\begin{eqnarray*} %
 D &=&(1+r^2)^3(1+3r^2)^{3/2}P(2-P+r^2)(P(1+r^2)+r^2(3-r^2))\\ %
 &\times&
 (P-1+2r^2)^{5/2}(P(1+r^2)-1-r^4)^{1/2}(r^2(P+1-r^2)+1)^{3/2}. %
 \end{eqnarray*}  %

It is easily seen that $D$ is non-negative.  Hence, to show that
$h(r)$ decreases monotonically, it suffices to show that $g=g(r):=
(c_0 + c_1P)+ (d_0 + d_1P)Q$ is non-negative for $0<r<1$.

We will show in Lemma \ref{Lemma Q} below that $0 < Q < 1$ for $0
< r < 1$.  Hence, to show that $g(r)$ is non-negative on $0<r<1$,
it will suffice, in view of the linearity of $g$ in $Q$, to show that
$$  %
g_0 = (c_0 + c_1P)+ (d_0 + d_1P) \cdot 0 \quad {\mbox{and}}  \quad
g_1 = (c_0 + c_1P)+ (d_0 + d_1P) \cdot 1  %
$$  %
are non-negative for $0<r<1$.  By Lemma \ref{Lemma P} below, we
have $s < P < t$, where %
\be  \label{equation 4.11}%
s =
\frac{19}{50}r^3-\frac{9}{10}r^2+\frac{1}{125}r+\frac{2827}{2000}
\quad {\mbox{and}} \quad  t =
\frac{9}{25}r^3-\frac{39}{50}r^2+\frac{1}{100}r+\frac{2829}{2000}. %
\ee %
 Hence, since $g_0$ and $g_1$ are linear in $P$, then we have %
$$  %
\min\{c_0+c_1s,c_0+c_1t\}\le g_0\le \max\{c_0+c_1s,c_0+c_1t\}  %
$$  %
and %
$$  %
\min\{c_0+d_0+(c_1+d_1)s,c_0+d_0+(c_1+d_1)t\}\le g_1\le
\max\{c_0+d_0+(c_1+d_1)s,c_0+d_0+(c_1+d_1)t\}.  %
$$  %



Since all the comparison expressions in these formulas are
polynomials in $r$ with rational coefficients, we can apply
 a Sturm sequence argument, see Chapter 5 of \cite{J}.
 This easily implies that $c_0+c_1s$, $c_0+c_1t$, $c_0+d_0+(c_1+d_1)s$, and $c_0+d_0+(c_1+d_1)t$
 % $ {g_0}_{\bigl |_{Z=s}}, \ {g_0}_{\bigl |_{Z=t}}, \ {g_1}_{\bigl |_{Z=s}}, \ {g_1}_{\bigl |_{Z=t}}  $
 are all non-negative for $0 \le r \le 1$. The proof is complete.  \hfill  $\Box$  %


\bl \label{Lemma P} %
Let $P$ be defined by (\ref{equation P and Q}) and let $s$ and $t$ be
defined by (\ref{equation 4.11}). Then, $s < P < t$ for $0 < r < 1$. %
\el  %

\textit{Proof.}  It suffices to show that $s^2 < P^2 < t^2$ for $0
< r < 1$.  We have %
$$   %\begin{eqnarray*}
P^2 - s^2 =  r^4 - 2r^2 + 2-
\left(\frac{19}{50}r^3-\frac{9}{10}r^2+\frac{1}{125}r+\frac{2827}{2000}\right)^2  %
$$  %
and %
$$  %
 t^2 - P^2 =
\left(\frac{9}{25}r^3-\frac{39}{50}r^2+\frac{1}{100}r+\frac{2829}{2000}\right)^2-r^4+2r^2-2. %
$$   %\end{eqnarray*}
Using a Sturm sequence argument, we can easily see that both $P^2
- s^2$ and $t^2 - P^2$ are non-negative for $0 < r < 1$. \hfill
$\Box$



\bl \label{Lemma Q} Let $Q$ be defined by (\ref{equation P and Q}).
Then, $ 0 < Q < 1 $ for $0 < r < 1$. \el

\textit{Proof.} It suffices to show that $0 <Q_1< e-1$ for $0 < r
< 1$,  where $Q_1 = \exp(Q)-1$.  We will show, in fact, that $0 <
Q_1 < 3/2$, which is equivalent to showing that $Q_1/(3/2-Q_1)
> 0$. We can write %
\be  \label{equation Q1}  %
\frac{Q_1}{3/2-Q_1} =
\frac{(4+4r^2)P-(4+4r^4)}{(9r^2+10+7r^4)-(7+7r^2)P}.  %
\ee  %
It is easily seen from (\ref{equation P and Q}) that $1 < P <
\sqrt{2}$ for $0 < r < 1$. Hence, it is clear that the numerator
in (\ref{equation Q1}) is positive. On the other hand, we have %
$$  %
(9r^2+10+7r^4)^2-(7+7r^2)^2P^2 = 270r^4+82r^2+126r^6+2
> 0, %
$$
which shows that the denominator in (\ref{equation Q1}) is
positive as well. The lemma is proved. \hfill $\Box$




\medskip

All the results established so far were used to prove
Theorem~\ref{Theorem 1}. Now we prove monotonicity properties of
the functions $\tau=\tau(r)$, $\beta=\beta(r)$, and $a=a(r)$.
Although not needed for our main proof they provide some
additional information about extremal configurations. The graph of
$a(r)$ is displayed in Figure~\ref{Figure_h_and_a} and the graphs
of functions $\tau(r)$ and $\beta(r)$  are displayed in Figure~\ref{Figure_tau_and_beta}.%


\begin{figure}
$$
\includegraphics[scale=.35,angle=-90]{graph-of-tau-alt}
\includegraphics[scale=.35,angle=-90]{graph-of-beta-alt}
$$
\caption{Functions $\tau(r)$ and $\beta(r)$.}
\label{Figure_tau_and_beta}
\end{figure}

\medskip

\bl \label{Lemma tau}  %
 The function $\tau=\tau(r)$ defined by (\ref{equation Thm1-C}) strictly decreases
 from $\sqrt{\sqrt 2-1}$ to $0$ as $r$ runs from $0$ to $1$. %
\el %

\textit{Proof.} It suffices to work with $\tau^2=\tau^2(r)$.
Differentiating $\tau^2$ we obtain %
$$  %
 \frac{d\tau^2}{dr} = \frac{2r}{P(1+3r^2)^2}\, p(r),  %
$$  %
where  $P$ is defined by (\ref{equation P and Q}) and %
$$   %
p(r) = -5+r^2-r^4+3r^6+(3-2r^2-3r^4)P.  %
$$  %
Hence, it suffices to show that $p(r)$ is negative for $0 \le r
\le 1$. It is easily seen that $P$ decreases from $\sqrt2$ to $1$ 
as $r$ varies from $0$ to $1$.  Hence, for  $0\le r\le 1$, we have
$1 \le P < 3/2$.  Suppose that 
\begin{eqnarray*}
c_1(r)&=&-5+r^2-r^4+3r^6 + (3-2r^2-3r^4)\cdot 1,\\
c_2(r)&=&-5+r^2-r^4+3r^6 + (3-2r^2-3r^4)\cdot (3/2).
\end{eqnarray*}
The linearity of $p$ with respect to $P$ implies that %
$$  %
 \min \{c_1(r),c_2(r)\} \le p(r) \le \max \{c_1(r),c_2(r)\}.  %
$$  %
Using a Sturm sequence argument, it is easily seen that both
$c_1(r)$ and $c_2(r)$ are negative for $0 \le r \le 1$. Thus,
$\tau^2(r)$ decreases on $0\le r\le 1$ and the lemma follows.
\hfill $\Box$


\bl  \label{Lemma beta}%
 Let $\beta=\beta(r)$ be defined by (\ref{equation Thm1-D}) with
 $\tau=\tau(r)$ defined by (\ref{equation Thm1-C}).  Then,
$\beta$ strictly increases from $0$ to $1$ as $r$ runs from $0$ to
$1$. %
\el  %

\textit{Proof.}  %
It suffices to show that $\beta^2$ is an increasing function of
$r$, which maps $[0,1]$ onto $[0,1]$.  We obtain, after some
algebra,  %
$$  %
\beta^2 = \frac{16 r^2 (2 r^4 + r^2 - 1) + 16 r^2 (1 + r^2) P}
{(1+r^2)^4 (1  + 2 r^2 - r^6) + (1+r^2)^4 (r^4+r^2)P},  %
$$  %
where $P$ is defined by (\ref{equation P and Q}). Differentiating
$\beta^2$, we find %
$$  %
 \frac{d\beta^2}{dr} = \frac{32 r (1-r^2)^2 (1+r^2)^5 }
{P((1+r^2)^4 (1  + 2 r^2 - r^6) + (1+r^2)^4 (r^4+r^2)P)^2}\,
p(r), %
$$  %
where %
$$  %
p(r) = -4 r^6 - r^4 - 5 r^2 + 2 + (4 r^4 + 5 r^2-1)P.  %
$$  %
 Hence,
it suffices to show that $p(r)$ is non-negative for $0 \le r \le 1$.
Now using a Sturm sequence argument, one can finish the proof
as in the previous lemma.
 \hfill  $\Box$  %



\medskip

Since $r=r(h)$ is monotonic on $0<h<4$, the parameters $\tau$ and
$\beta$ in the definition of the extremal function $f_h$ of
Theorem~\ref{Theorem 1} are monotonic functions of $h$. It is
worth mentioning that the third natural parameter, $a=a(h)$, which
gives the length of the vertical segment of the non-free boundary,
\emph{is not} monotonic in $h$. It is easy to see that the disk
$\{w:\,|w-1|\le 1\}$ and segment $[0,4]$ are the limit extremal
configurations for the problem under consideration. Thus, $a=0$ in
both limit cases. Our next lemma shows however that $a=a(r)$
considered as a function of $r$ has only one local maximum on
$0<r<1$.

\bl  \label{Lemma a}  %
Let $a=a(r)$ be defined by (\ref{equation a}) with $\tau$ and
$\beta$ defined by (\ref{equation Thm1-C}) and (\ref{equation Thm1-D}).
Then, there is a unique $r_1$, $0<r_1<1$, such that $a(r)$
strictly increases as $r$ varies from $0$ to $r_1$  and strictly
decreases as $r$ varies from $r_1$ to $1$.  %
\el  %

\textit{Proof.}
Upon integration, $a(r)$ can be expressed as an explicit function of $r$ which is
a combination of polynomials, square
roots, and arctangents. We give here an argument that is reminiscent of the
argument given in the proof of Lemma \ref{Lemma h}, omitting some of the technical details.
For convenience, we set $r_0 = 53/100$ and $r_2 = 57/100$.

Differentiating $a(r)$ with respect to $r$ we
obtain a representation
$$a'(r) = 4r(1-r^2)\frac{(c_0 + c_1\,P)+(d_0 + d_1\,P)G(r)}{D_1(r)}$$
where $P = (2-2r^2+r^4)^{1/2}$, the functions $G$ and $D_1$ are
non-negative on $(0,1)$ and $c_0, \ c_1, \ d_0$ and $d_1$ are
polynomials in $r$ with rational coefficients. We will show that
there exists an $r_1$ such that $a'(r) > 0$ on $(0,r_1)$ and $a'(r) <
0$ on $(r_1,1)$.

Using the linearity of the terms
$c_0 + c_1\,P$ and $d_0 + d_1\,P$ in $P$ and the estimates on $P$ given in
Lemma \ref{Lemma P}, one can give a Sturm sequence argument to show that
$c_0 + c_1\,P > 0$ and $d_0 + d_1\,P > 0$ on the interval $(0,r_0)$ and that
$c_0 + c_1\,P < 0$ and $d_0 + d_1\,P < 0$ on the interval $(r_2,1)$.  Hence,
$a'(r) > 0$ on $(0,r_0)$ and $a'(r) < 0$ on $(r_2,1)$.

We define $n(r) = (c_0 + c_1\,P)+(d_0 + d_1\,P)G(r)$.  Differentiating $n(r)$ with respect to $r$
we obtain a representation
$$n'(r) = 2r\tau^2\frac{({\tilde c}_0 + {\tilde c}_1\,P)+({\tilde d}_0 + {\tilde d}_1\,P)G(r)}{D_2(r)}$$
where the function $D_2$ is non-negative on $(0,1)$, $\tau$ is
defined by (\ref{equation Thm1-C}) and ${\tilde c}_0, \ {\tilde
c}_1, \ {\tilde d}_0$ and ${\tilde d}_1$ are polynomials in $r$
with rational coefficients.

Using the linearity of terms
${\tilde c}_0 + {\tilde c}_1\,P$ and ${\tilde d}_0 + {\tilde d}_1\,P$ in $P$ and
the estimates on $P$ given in
Lemma \ref{Lemma P}, one can give a Sturm sequence argument to show that
${\tilde c}_0 + {\tilde c}_1\,P < 0$ and ${\tilde d}_0 + {\tilde d}_1\,P < 0$ on the
interval $(r_0,r_2)$ and, hence, that $n(r)$ is strictly decreasing on the interval $(r_0,r_2)$ and
changes sign exactly once. Consequently, $a'(r)$ changes sign exactly once on $(r_0,r_2)$.

The value $r_1$ is the unique solution of $n(r) = 0$, which lies in the interval $(r_0,r_2)$.
\hfill  $\Box$  %


\section{Some remarks and problems}  %
\setcounter{equation}{0} %

\textbf{(a) Omitted area problem.} The following problem proposed
by A.~W.~Goodman \cite{G} can be considered as a prototype of all
omitted area problems with geometrical constraints:
%Our original interest to problems with some knowledge on a certain portion of a
%compact set came from the following problem proposed by A.~W.~Goodman:
Find $A:=\inf_{f\in S} \,\{ {\rm{Area}}\,(f(\mathbb{D})\cap
\mathbb{D})\}$ over the standard class $S$ of univalent functions
$f$ in $\mathbb{D}$ with $f(0)=0$, $f'(0)=1$.

To our knowledge, this problem remains open although many
important properties of extremal functions have been proved since
1949. Here we summarize some of them. If $f\in S$, $f$ is extremal for
$A$ and $f(1)=\infty$, then $D=f(\mathbb{D})$ is circularly
symmetric w.r.t. $\mathbb{R}_0$ %(up to rotation about the origin)
and there exist $\theta_1,\theta_2$, and $\beta$ such that
 $0<\theta_1<\theta_2<\pi$, $0<\beta<1$,
and   $f$ satisfies the following boundary conditions: %
\begin{enumerate} %
\item[(a)] $\Im\,f(e^{i\theta})=0$ for $0<|\theta|\le \theta_1$; %
\item[(b)] $|f(e^{i\theta})|=1$ for $\theta_1<|\theta|<\theta_2$; %
\item[(c)] $|f'(e^{i\theta})|=\beta$ for
$\theta_2<\theta<2\pi-\theta_2$; %
\item[(d)] $f'$ has a non-zero continuous extension to
$\mathbb{D}\cup \{e^{i\theta}:\,\theta_1<\theta<2\pi-\theta_1\}$
which is H\"{o}lder-continuous with exponent $1/2$; %
\item[(e)] $|f'(e^{i\theta})|$ strictly decreases in
$\theta_1<\theta<\theta_2$; %
\item[(f)] there is a $\theta_0$, $0<\theta_0<\theta_1$ such that
$|f'(e^{i\theta})|$ strictly decreases from $+\infty$ to $\beta_1$, where 
$\beta_1>\beta$, and strictly increases from $\beta_1$ to $+\infty$
in $0<\theta<\theta_0$ and $\theta_0<\theta<\theta_1$,
respectively. %
\end{enumerate}  %

Observations (a) and (b) were made by Barnard and Suffridge, see
\cite[p. 536]{BrC}. Condition (d) was proved by J.~Lewis \cite{L}
who also proved that (c) holds true for all $\theta$ except the
set $I=\{e^{i\theta}:\,\Im\,f(e^{i\theta})=0\}$ which may consists
of at most a finite number of closed arcs. The inequality
$\beta<1$ and conditions (e), (f), and (c) without the above
mentioned exception were established in \cite{BS}.

\smallskip

The conclusion of Lemma~\ref{Lemma 2.2}\textbf{(d)}
that the vertical non-free boundary is not degenerate, i.e., that
there is a strict inequality $\theta_1<\theta_2$  for
the parameters $\theta_1$ and $\theta_2$ of this iceberg-type problem is reminscient of the
conclusion in \cite{L} that there is a strict inequality $\theta_1<\theta_2$  for
the parameters $\theta_1$ and $\theta_2$ of Goodman's omitted
area problem.  With minor modifications, the proof in \cite{L} that
$\theta_1<\theta_2$ for Goodman's omitted area problem could have been modified to
prove Lemma~\ref{Lemma 2.2}\textbf{(d)}.  In this paper, we have given an independent
proof of Lemma~\ref{Lemma 2.2}\textbf{(d)} and we mention here that, alternatively,
with minor modifications the proof of Lemma~\ref{Lemma 2.2}\textbf{(d)} could be used
to show that $\theta_1<\theta_2$ for Goodman's omitted area problem as well.  
Approximations to the exact value of $A$ have been given in \cite{BP,BT} by 
different numerical methods.  In particular, \cite{BT} suggests that 
$A = 0.2385813284\pi$, where all explicitly shown digits are exact.

\medskip

\textbf{(b) Width of the invisible part of the iceberg.} The
method of this paper can be also applied to find the extremal
function for Problem~(\ref{equation maximals})\textbf{(b)} if one
can show a priori that the free boundary of the extremal is smooth enough. One
difference compared to Problem~(\ref{equation
maximals})\textbf{(a)} is that the extremal configurations now
\emph{do not} possess circular symmetry although they still
possess Steiner symmetry. In view of this lack of symmetry, we cannot 
apply the local variations developed in
Section~\ref{Section 2} since the boundary may be non-rectifiable.
Perhaps, the necessary smoothness can be achieved by applying a
more powerful technique such as that of J.~Lewis \cite{L}
mentioned in the Introduction.

\medskip

\textbf{(c) Safe distance from the iceberg.} The situation with
Problem~(\ref{equation maximals})\textbf{(c)} differs from the
other two cases. To explain this, we start with the limiting case
when the whole iceberg is observable, i.e. when $\Area (E_+)=\pi$.
Then, of course, $E$ coincides with the disk $\{w:\,|w-(1+i)|\le
1\}$  up to translation along the real axis.

This disk has a contact point with the surface of interface at
$z=1$ and a contact point with the front line, which coincides
with the imaginary axis, at $z=i$. These two contact points
represent the non-free boundary in this limiting case. It is
reasonable to expect that for icebergs with visible area slightly
less than $\pi$, the extremal configurations will have two
disjoint segments, vertical and horizontal, as their non-free
boundary. If so, then transplanting the problem into the auxiliary
$s$-plane as in Section~\ref{Section 3}, we have to deal with the
omitted area problem for functions defined in a doubly-connected
domain. To our knowledge, there are no known solutions of 
problems of this kind.

\medskip

\textbf{(d) Convex icebergs.} Let ${\mathcal{F}}^c$ denote the
collection of all convex compact sets $E$ in ${\mathcal{F}}$. It
will be interesting to study problems (\ref{equation maximals})
for the class ${\mathcal{F}}^c$. Since there are more available
methods for convex sets and functions, there is a chance that
known techniques may give complete solutions to all three
problems.


\begin{thebibliography}{6.8in}

\bibitem{AShS1}
D.~Aharonov, H.~S.~Shapiro, and A.~Yu.~Solynin, 
\textit{A minimal area problem in conformal mapping}. J. Analyse Math.
\textbf{78} (1999), 157--176.

\bibitem{AShS2}
D.~Aharonov, H.~S.~Shapiro, and A.~Yu.~Solynin, 
\textit{Minimal area problems for functions with integral representation}. J. Analyse Math. 
\textbf{98} (2006), 83--111.

\bibitem{BT} L. Banjai, L. Trefethen,
\textit{Numerical solution of the omitted area problem of univalent function theory.} Comp. Methods Func. Theory
\textbf{1} (2001), no. 1, 259--273.

\bibitem{BL} R. W. Barnard, J. L. Lewis,
\textit{On the omitted area problem.} Michigan Math. J.
\textbf{34} (1987), 13--22.

\bibitem{BP} R. W. Barnard, K. Pearce,
\textit{Rounding Corners of Gearlike Domains and the Omitted Area Problem.} J. Comp. Appl. Math.
\textbf{14} (1986), 217--226; 
\textit{Numerical Conformal Mapping.} North Holland, 1986.

\bibitem{BPS1} R. W. Barnard, K. Pearce, and A. Yu. Solynin,
\textit{An isoperimetric inequality for logarithmic capacity.} Annales Academi\ae \ Scientiarum Fennic\ae. Mathematica
\textbf{27} (2002), 419--436.

\bibitem{BRS1} R. W. Barnard, C. Richardson, and A. Yu. Solynin,
\textit{Concentration of area in half-planes.} Proc. Amer. Math. Soc. 
\textbf{133} (2005), no. 7, 2091--2099.

\bibitem{BRS2} R. W. Barnard, C. Richardson, and A. Yu. Solynin,
\textit{A minimal area problem for nonvanishing functions.} Algebra i Analiz 
\textbf{18} (2006), no. 1, 35-54; 
English translation in: St. Petersburg Math. J. 
\textbf{18} (2007), no. 1, 21--36.

\bibitem{BS} R. W. Barnard and A. Yu. Solynin, 
\textit{Local variations and minimal area problem for Carath\'{e}odory functions.} Indiana U. Math. J. 
\textbf{53} (2004), no. 1, 135--167.

\bibitem{BrC} D. Brannan and J.~Clunie,
\textit{Aspects of Contemporary Complex Analysis.} Academic Press, New York, 1980.

\bibitem{D} V. N. Dubinin, 
\textit{Symmetrization in geometric theory of functions of a complex variable.}  Uspehi Mat. Nauk
\textbf{49} (1994), 3-76 (in Russian); 
English translation in: Russian Math. Surveys 
\textbf{49}: 1 (1994), 1--79.

\bibitem{Du} P. Duren, 
\textit{Univalent Functions.} Springer-Verlag, 1992.

\bibitem{G} A. Goodman,
\textit{Note on regions omitted by univalent functions.} Bull. Amer. Math. Soc. 
\textbf{55} (1949), 363--369.

\bibitem{H} W. K. Hayman, 
\textit{Multivalent Functions.} Second Edition. Cambridge Tracts in Mathematics, 110. Cambridge Univ. Press, Cambridge, 1994.

\bibitem{J} N. Jacobson, 
\textit{Basic Algebra. I.} Second Edition.  W. H. Freeman and Company, New York, 1985. % MR0780184 (86d:00001)

\bibitem{L}  J. L. Lewis,
\textit{On the minimal area problem.} Indiana Univ. Math. J.
\textbf{34} (1985), 631--661.

\bibitem{P} Ch. Pommerenke, 
\textit{Boundary Behaviour of Conformal Maps.} Springer-Verlag, 1992.
ath. J. \textbf{8} (1997), 1015--1038.
%
\bibitem{S} A. Yu. Solynin,
\textit{A Functional inequalities via polarization.} Algebra i Analiz 
\textbf{8} (1996), 145-185; 
English translation in: St. Petersburg Math. J. 
\textbf{8} (1997), 1015--1038.
\end{thebibliography}

\end{document}
